\ssr{ПРИЛОЖЕНИЕ А}

\begin{lstlisting}[language=C++, caption={Первая итерация с помощью DeepSeek}, captionpos=b, label=list:deepseek-1]
import os
import re
import sys
from typing import List, Tuple, Optional
import PyPDF2


def find_non_cyrillic_in_tokens(pdf_path: str) -> Tuple[bool, List[Tuple[str, Tuple[int, int]]]]:
    pattern = re.compile(
        r'''
        (?:^|\s)                   
        ([а-яА-ЯёЁ0-9\-.,:;!?()"']+?[^а-яА-ЯёЁ0-9\s\-.,:;!?()"'][а-яА-ЯёЁ0-9\-.,:;!?()"']*?) 
        (?=\s|$)                   
        ''', 
        re.VERBOSE | re.UNICODE
    )
    
    results = []
    
    try:
        with open(pdf_path, 'rb') as file:
            pdf_reader = PyPDF2.PdfReader(file)
            
            for page_num in range(len(pdf_reader.pages)):
                page = pdf_reader.pages[page_num]
                page_text = page.extract_text()
                
                if page_text:
                    lines = page_text.split('\n')
                    
                    for line_num, line in enumerate(lines, 1):
                        matches = pattern.findall(line)
                        
                        for match in matches:
                            if match:
                                results.append((match, (page_num + 1, line_num)))
    
    except Exception as e:
        print(f"Ошибка при обработке файла {pdf_path}: {e}")
        return False, []
    
    found = len(results) > 0
    return found, results


def process_pdf_files_in_directory(directory_path: str = "files") -> dict:
    if not os.path.exists(directory_path):
        print(f"Директория '{directory_path}' не найдена.")
        return {}
    
    pdf_files = []
    for file in os.listdir(directory_path):
        if file.lower().endswith('.pdf'):
            pdf_files.append(os.path.join(directory_path, file))
    
    if not pdf_files:
        print(f"В директории '{directory_path}' не найдено PDF-файлов.")
        return {}
    
    print(f"Найдено {len(pdf_files)} PDF-файлов для обработки.")
    
    results = {}
    for pdf_file in pdf_files:
        print(f"\nОбработка файла: {os.path.basename(pdf_file)}")
        found, matches = find_non_cyrillic_in_tokens(pdf_file)
        results[pdf_file] = {
            'found': found,
            'matches': matches,
            'filename': os.path.basename(pdf_file)
        }
        
        if found:
            print(f"Найдены токены с некириллическими символами: {len(matches)} шт.")
            for match, (page, line) in matches:
                print(f"- '{match}' на странице {page}, строке {line}")
        else:
            print("Некириллических символов в кириллических токенах не найдено.")
    
    return results


def main():
    if len(sys.argv) < 2:
        print("Использование:")
        print("  python script.py <путь_к_pdf_файлу>")
        print("  python script.py <путь_к_директории>")
        print("  python script.py --dir     (для обработки директории 'files')")
        print("\nПримеры:")
        print("  python script.py document.pdf")
        print("  python script.py /путь/к/файлу/document.pdf")
        print("  python script.py /путь/к/директории")
        print("  python script.py --dir")
        return
    
    input_path = sys.argv[1]
    
    if input_path == "--dir":
        process_pdf_files_in_directory()
    elif os.path.isdir(input_path):
        process_pdf_files_in_directory(input_path)
    elif os.path.isfile(input_path) and input_path.lower().endswith('.pdf'):
        print(f"Обработка файла: {input_path}")
        found, matches = find_non_cyrillic_in_tokens(input_path)
        
        print(f"\nРезультаты для файла: {os.path.basename(input_path)}")
        print(f"Найдены токены с некириллическими символами: {found}")
        
        if found:
            print(f"Количество найденных токенов: {len(matches)}")
            for i, (match, (page, line)) in enumerate(matches, 1):
                print(f"{i}. '{match}' - страница {page}, строка {line}")
        else:
            print("Токенов с некириллическими символами не найдено.")
    else:
        print(f"Ошибка: '{input_path}' не является PDF-файлом или директорией.")


if __name__ == "__main__":
    main()
\end{lstlisting}

\begin{lstlisting}[language=C++, caption={Вторая итерация с помощью DeepSeek}, captionpos=b, label=list:deepseek-2]
import os
import re
import sys
from typing import List, Tuple
import PyPDF2


def find_non_cyrillic_in_tokens(pdf_path: str) -> Tuple[bool, List[Tuple[str, Tuple[int, int]]]]:
    pattern = re.compile(
        r'''
        (?:^|\s)               
        (
          (?:                      
            [а-яА-ЯёЁ]+            
            |                      
            [а-яА-ЯёЁ][0-9\-.,:;!?()"'%]* 
          )
          [^а-яА-ЯёЁ0-9\s\-.,:;!?()"'%] 
          .*?                       
        )
        (?=\s|$)
        ''', 
        re.VERBOSE | re.UNICODE
    )
    
    pattern2 = re.compile(
        r'''
        (?:^|\s)
        (                          
          [а-яА-ЯёЁ]               
          [а-яА-ЯёЁ0-9\-.,:;!?()"'%]* 
          [^а-яА-ЯёЁ0-9\s\-.,:;!?()"'%]
          [а-яА-ЯёЁ0-9\-.,:;!?()"'%]* 
        )
        (?=\s|$)
        ''',
        re.VERBOSE | re.UNICODE
    )

    simple_pattern = re.compile(
        r'''
        (?:^|\s)                  
        (                  
          [^\s]+     
        )
        (?=\s|$)            
        ''',
        re.VERBOSE | re.UNICODE
    )
    
    results = []
    
    try:
        with open(pdf_path, 'rb') as file:
            pdf_reader = PyPDF2.PdfReader(file)
            
            for page_num in range(len(pdf_reader.pages)):
                page = pdf_reader.pages[page_num]
                page_text = page.extract_text()
                
                if page_text:
                    lines = page_text.split('\n')
                    
                    for line_num, line in enumerate(lines, 1):
                        potential_tokens = simple_pattern.findall(line)
                        
                        for token in potential_tokens:
                            has_cyrillic = bool(re.search(r'[а-яА-ЯёЁ]', token))
                            
                            if has_cyrillic:
                                has_non_cyrillic = bool(
                                    re.search(
                                        r'[^а-яА-ЯёЁ0-9\s\-.,:;!?()"'']',
                                        token
                                    )
                                )
                                
                                if has_non_cyrillic:
                                    cyrillic_chars = re.findall(r'[а-яА-ЯёЁ]', token)
                                    total_chars = len(token)
                                    
                                    if len(cyrillic_chars) / total_chars >= 0.3:
                                        results.append((token, (page_num + 1, line_num)))
    
    except Exception as e:
        print(f"Ошибка при обработке файла {pdf_path}: {e}")
        return False, []
    
    found = len(results) > 0
    return found, results


def process_pdf_files_in_directory(directory_path: str = "files") -> None:
    print("Программа для поиска некириллических символов в кириллических токенах")
    print("=" * 60)
    
    if not os.path.exists(directory_path):
        print(f"Ошибка: Директория '{directory_path}' не найдена.")
        print("Создайте директорию 'files' и поместите в нее PDF-файлы для обработки.")
        return
    
    pdf_files = []
    for file in os.listdir(directory_path):
        if file.lower().endswith('.pdf'):
            pdf_files.append(os.path.join(directory_path, file))
    
    if not pdf_files:
        print(f"В директории '{directory_path}' не найдено PDF-файлов.")
        print("Поместите PDF-файлы в директорию 'files'.")
        return
    
    print(f"Найдено {len(pdf_files)} PDF-файлов для обработки:")
    for i, file in enumerate(pdf_files, 1):
        print(f"  {i}. {os.path.basename(file)}")
    print("=" * 60)
    
    all_results = []
    for pdf_file in pdf_files:
        filename = os.path.basename(pdf_file)
        print(f"\nОбработка файла: {filename}")
        print("-" * 40)
        
        found, matches = find_non_cyrillic_in_tokens(pdf_file)
        
        if found:
            print(f"Найдены кириллические токены: {len(matches)} шт.")
            for match, (page, line) in matches:
                print(f"  - '{match}' (страница {page}, строка {line})")
            all_results.append((filename, matches))
        else:
            print("Токенов не найдено.")
    
    print("\n" + "=" * 60)
    print("ИТОГОВЫЙ ОТЧЕТ")
    print("=" * 60)
    
    if all_results:
        total_tokens = sum(len(matches) for _, matches in all_results)
        print(f"Всего найдено кириллических токенов с некириллическими символами: {total_tokens}")
        print("Файлы с найденными токенами:")
        for filename, matches in all_results:
            print(f"  - {filename}: {len(matches)} токенов")
    else:
        print("Во всех файлах отсутствуют.")
    
    print("\nОбработка завершена.")


def main():
    process_pdf_files_in_directory()


if __name__ == "__main__":
    main()
\end{lstlisting}



\begin{lstlisting}[language=C++, caption={Первая итерация запроса в ChatGPT 5.1}, captionpos=b, label=list:chatgpt-1]
    import os
from PyPDF2 import PdfReader


def is_cyrillic_letter(ch):
    return 'А' <= ch <= 'Я' or 'а' <= ch <= 'я' or ch == 'Ё' or ch == 'ё'


def check_pdf_for_invalid_cyrillic_tokens(pdf_path):
    reader = PdfReader(pdf_path)
    results = []

    for page_number, page in enumerate(reader.pages, start=1):
        text = page.extract_text()
        if not text:
            continue

        text = text.replace('\u00A0', ' ')
        lines = text.splitlines()

        for line_number, line in enumerate(lines, start=1):
            tokens = line.split()
            for token in tokens:
                has_cyrillic = False
                has_foreign_letter = False

                for ch in token:
                    if ch.isalpha():
                        if is_cyrillic_letter(ch):
                            has_cyrillic = True
                        else:
                            has_foreign_letter = True

                if has_cyrillic and has_foreign_letter:
                    results.append((token, page_number, line_number))

    if results:
        return True, results
    return False, []


def process_all_pdfs_in_directory(directory):
    reports = {}
    for root, _, files in os.walk(directory):
        for file in files:
            if file.lower().endswith(".pdf"):
                path = os.path.join(root, file)
                reports[path] = check_pdf_for_invalid_cyrillic_tokens(path)
    return reports


def print_report(reports):
    for pdf_path, (exists, entries) in reports.items():
        print("=" * 60)
        print(f"Файл: {os.path.basename(pdf_path)}")
        if not exists:
            print("Нарушения не обнаружены")
        else:
            print("Обнаружены некорректные кириллические токены:")
            for token, page, line in entries:
                print(f"  Страница {page}, строка {line}: {token}")
        print()


if __name__ == "__main__":
    reports = process_all_pdfs_in_directory("files")
    print_report(reports)
report/08-appendix.tex
\end{lstlisting}

\begin{lstlisting}[language=C++, caption={Первая итерация запроса в Qwen-Coder}, captionpos=b, label=list:qwen-1]
import os
from PyPDF2 import PdfReader


def is_cyrillic_letter(ch):
    return 'А' <= ch <= 'Я' or 'а' <= ch <= 'я' or ch == 'Ё' or ch == 'ё'


def check_pdf_for_invalid_cyrillic_tokens(pdf_path):
    reader = PdfReader(pdf_path)
    results = []

    for page_number, page in enumerate(reader.pages, start=1):
        text = page.extract_text()
        if not text:
            continue

        text = text.replace('\u00A0', ' ')
        lines = text.splitlines()

        for line_number, line in enumerate(lines, start=1):
            tokens = line.split()
            for token in tokens:
                has_cyrillic = False
                has_foreign_letter = False

                for ch in token:
                    if ch.isalpha():
                        if is_cyrillic_letter(ch):
                            has_cyrillic = True
                        else:
                            has_foreign_letter = True

                if has_cyrillic and has_foreign_letter:
                    results.append((token, page_number, line_number))

    if results:
        return True, results
    return False, []


def process_all_pdfs_in_directory(directory):
    reports = {}
    for root, _, files in os.walk(directory):
        for file in files:
            if file.lower().endswith(".pdf"):
                path = os.path.join(root, file)
                reports[path] = check_pdf_for_invalid_cyrillic_tokens(path)
    return reports


def print_report(reports):
    for pdf_path, (exists, entries) in reports.items():
        print("=" * 60)
        print(f"Файл: {os.path.basename(pdf_path)}")
        if not exists:
            print("Нарушения не обнаружены")
        else:
            print("Обнаружены некорректные кириллические токены:")
            for token, page, line in entries:
                print(f"  Страница {page}, строка {line}: {token}")
        print()


if __name__ == "__main__":
    reports = process_all_pdfs_in_directory("files")
    print_report(reports)
\end{lstlisting}



\begin{lstlisting}[language=C++, caption={Вторая итерация запроса в Qwen-Coder}, captionpos=b, label=list:qwen-2]
import os
import re
from pathlib import Path
from PyPDF2 import PdfReader
from PyPDF2.errors import PdfReadError


def find_non_cyrillic_in_cyrillic_tokens(pdf_path):
    if not os.path.isfile(pdf_path):
        return (False, [])
    
    try:
        reader = PdfReader(pdf_path)
    except Exception:
        return (False, [])
    
    results = []
    # Паттерн для извлечения потенциальных "слов":
    # буквы (любые), цифры, дефис, апострофы (включая Unicode ’)
    word_pattern = re.compile(r"[a-zA-Zа-яА-ЯёЁ0-9\-\'\u2019]+")

    for page_num, page in enumerate(reader.pages, start=1):
        try:
            text = page.extract_text()
        except Exception:
            continue
        if not text:
            continue
        lines = text.split('\n')
        for line_num, line in enumerate(lines, start=1):
            tokens = word_pattern.findall(line)
            for token in tokens:
                if re.search(r'[а-яА-ЯёЁ]', token):  # есть кириллица?
                    # Проверяем: есть ли символы вне разрешённого кириллического набора?
                    if re.search(r'[^а-яА-ЯёЁ0-9\-\'\u2019]', token):
                        results.append((token, (page_num, line_num)))
                        break
    return (len(results) > 0, results)


def process_all_pdfs_in_directory(directory="files"):
    results_dict = {}
    dir_path = Path(directory)
    if not dir_path.exists() or not dir_path.is_dir():
        return results_dict
    for pdf_file in sorted(dir_path.glob("*.pdf")):
        pdf_path_str = str(pdf_file)
        result = find_non_cyrillic_in_cyrillic_tokens(pdf_path_str)
        results_dict[pdf_path_str] = result
    return results_dict


def main():
    result = process_all_pdfs_in_directory()
    print("\n📊 РЕЗУЛЬТАТЫ АНАЛИЗА PDF-ФАЙЛОВ\n" + "="*50)
    if not result:
        print("📭 Нет PDF-файлов в директории 'files'.")
        return

    any_found = False
    for pdf_path, (found, matches) in result.items():
        print(f"\n📄 {os.path.basename(pdf_path)}")
        if found:
            any_found = True
            for token, (page, line) in matches:
                print(f"   • '{token}' — стр. {page}, строка {line}")
        else:
            print("   ✔ Некириллических символов в кириллических словах не найдено.")
    
    if not any_found:
        print("\n✅ Все файлы соответствуют требованиям.")
    else:
        print("\n❗ Обнаружены кириллические токены с некириллическими символами.")


if __name__ == "__main__":
    main()
\end{lstlisting}